% Archival extraction from draft/main.tex on 2026-10-10.
% The manuscript block below is preserved verbatim from the section beginning
% "Deferred ML protocol after the exact theorem gate" through the line before
% "Supporting context outside the reduction claim".
% It is not included by draft/main.tex and is not current experimental evidence.

\section{Deferred ML protocol after the exact theorem gate}
\label{sec:next-experiment}
This section preserves a secondary prospective protocol, not a report of new
training runs and not the paper's active computation. Its local question is
whether a fourth-order equation-derived prior improves the learned family
beyond the quadratic polynomial already needed to impose a third-order IVP's
data. A separate comparison tests the benefit of restricting the field to
verified invariant coordinates. These runs are admissible only as a separately
certified fixed-$m=3$ study after its reduction, data, flux and reference-solver
gates pass; they cannot answer the unrestricted near-cubic theorem.

\subsection{Mathematical and reference-solution prerequisites}
Use the nondimensional printed scalar benchmark at
$m=3$, $g=1$, $\beta_{\rm tr}=0.1$, $\kappa=0.1/3$.
Keep $\delta=m-3=0$. On $z\in[1,1.5]$ and $t\in[1,2]$, let
$K=\sqrt[3]{45}$ and vary only the data parameter $a\in[-0.02,0.02]$:
\begin{equation}
 F(1;a)=K,\qquad F'(1;a)=K/3,\qquad
 F''(1;a)=-2K(1+a)/9.
 \label{eq:experiment-ivp}
\end{equation}
The corresponding original data on $\Gamma_0=\{st^2=1\}$ are
\begin{equation}
 \phi=K,\qquad t^{-2}\phi_s=K/3,\qquad
 t^{-4}\phi_{ss}=-2K(1+a)/9.
 \label{eq:experiment-original-data}
\end{equation}
Their prolonged $X_5$ action vanishes, and their pullbacks give
equation~\eqref{eq:experiment-ivp}. This is a constructed compatible data
family, not a physically sourced canal boundary problem. At $a=0$ the exact
profile~\eqref{eq:sve-exact} supplies a check. Normal-form existence and
continuous parameter dependence give a neighborhood of a regular solution;
the entire proposed interval/domain still needs verified guard bounds or a
clearly labeled numerical regularity check.

Before training, regenerate the exact chart, pullback and data certificates
from available source. Recover the newer generic backend only if its
capabilities are claimed. For every reference profile record
$\min F'$, $\min|c_{\rm lead}|$, chart/domain exits and independent solver
agreement. Solve all 17 distinct parameter cases below with DOP853 at
$(\mathrm{rtol},\mathrm{atol})=(10^{-9},10^{-12})$,
$(10^{-11},10^{-14})$, and $(10^{-13},10^{-15})$, and cross-check with Radau.
For the proposed $10^{-5}$ profile-error target, require at most $10^{-8}$
reference discrepancy on the evaluation grid. If guards fail, split or
shrink the family before training and disclose the change. Finite mesh
checks are not a uniform interval certificate.

For an analytical certificate on this reduced IVP, use
Proposition~\ref{prop:ode-error-core} with a verified regular tube,
a bound on $D_YG$, the leading-factor margin and a continuous residual
bound. These are additional deliverables beyond solver agreement.

Write the regular ODE as $F'''=G(z,F,F',F'')$. With $h=z-1$, define
\begin{align}
 P_2(z,a)&=K+(K/3)h-K(1+a)h^2/9,\\
 P_4(z,a)&=P_2+c_3(a)h^3+c_4(a)h^4,\label{eq:experiment-priors}\\
 c_3&=G(1,Y_0)/6,\qquad
 c_4=\bigl(G_z+F'G_F+F''G_{F'}+G\,G_{F''}\bigr)(1,Y_0)/24,\nonumber
\end{align}
where $Y_0=(F(1),F'(1),F''(1))$. Verify this recurrence symbolically and
against the exact $a=0$ series. Choosing $P_4$ does not prove that its
truncation is accurate over the whole training interval.

\subsection{Ablations that actually isolate the analytical prior}
\begin{center}\small
\renewcommand{\arraystretch}{1.16}
\begin{tabularx}{\linewidth}{@{}p{3cm}p{5.1cm}X@{}}
\toprule
Arm & Trial function & Question isolated \\
\midrule
B: reduced hard-IVP PINN & $\widehat F=P_2+h^3N_\theta(z,a)$ & Baseline with exactly the same three initial data. \\
H: added analytical prior & $\widehat F=P_4+h^3N_\theta(z,a)$ & Extra equation-derived coefficients with the same remainder factor and neural architecture. \\
H5: extra hard series jets & $\widehat F=P_4+h^5N_\theta(z,a)$ & Additional exact Taylor constraints and different conditioning; a separate ablation. \\
D: direct PDE PINN & $\widehat\phi=P_2(st^2,a)+(st^2-1)^3N_\theta(t,s,a)$ & Same supplied data, without imposing $Q_5=0$ on the trainable field. \\
Controls & $P_2$, $P_4$, exact $a=0$ profile; converged classical solver & Accuracy and total cost without neural training. \\
\bottomrule
\end{tabularx}
\end{center}

A ``hard-IVP PINN without an analytical part'' is misleading here: imposing
three data jets already builds $P_2$ into the trial function. Moreover B and H
have the same unconstrained smooth profile class up to a shift of the neural
remainder by $c_3+c_4h$; any observed advantage concerns the chosen finite
network, initialization and optimization, not elimination of unknown
degrees of freedom by $P_4$. H5 adds stronger exact local information and
must not be used alone to attribute an improvement to the prior.
The direct arm knows the supplied surface $st^2=1$ through its data lifting;
it is not a baseline with no information about the data geometry.

For the fixed cubic benchmark, equation~\eqref{eq:sve-defect} implies
$\pr^{(3)}X_5\E_3=-\E_3$. Squaring that Lie derivative of the residual
merely reweights the original squared-residual loss. It provides no
independent symmetry information in this case. A graph penalty
$\|Q_5[\widehat\phi]\|^2$, or a genuinely different symmetry-informed
baseline, must be named separately from the residual-Lie-derivative
construction of \citet{sympinn}. Exact reduction imposes $Q_5=0$ by
construction, but restricts the solution family.

In the full adapted PDE, the coefficient of the highest normal derivative
$v_{zzz}$ on the initial surface is $c_{\rm lead}/t=2K/t\ne0$ for these
data. The equation is analytic on the regular branch. The
Cauchy--Kowalevski theorem therefore supplies local existence and uniqueness
for the analytic data, consistently with the reconstructed analytic ODE
solution. This is a local analytic statement. A common domain reaching
$z=1.5$, smooth-problem stability and an error bound for approximate fields
still need their own analysis. Limit the numerical comparison to the stated
regular reference branch until those stronger claims are justified.

\paragraph{Classical approximation and launch criteria.}
Alongside the classical IVP solver, include a tensor-product Chebyshev
approximation or spline interpolant in $(z,a)$, built from declared
reference data. Report its label-generation cost separately from the
label-free residual training of B/H. This comparator tests whether a neural
family surrogate is useful beyond ordinary smooth interpolation; it is
not a matched zero-label baseline. Also retain an adaptive Taylor or
polynomial approximation of each profile to separate neural benefit from
the benefit of the analytical expansion itself.

Training starts only after (i) the exact scalar model and branch are frozen,
(ii) symbolic reconstruction and data checks are reproducible,
(iii) the reference family passes the declared convergence and guard tests,
and (iv) the loss measure, validation split, primary metric and cost
accounting are fixed. These gates authorize an empirical scalar-benchmark
experiment. A certified experiment additionally requires interval-wide
solution-tube and continuous-residual bounds; a physical SVE experiment
additionally requires resolving the bed-transport derivation. Neither
requirement is supplied by a small sampled training loss.

\subsection{Splits, optimization budget and smoothness}
Use disjoint data-parameter sets, with the final test set reserved before
hyperparameter selection:
\begin{align*}
 \mathcal T_a&=\{-0.02,-0.015,-0.01,0,0.01,0.015,0.02\},\\
 \mathcal V_a&=\{-0.0125,-0.005,0.005,0.0125\},\\
 \mathcal S_a&=\{-0.018,-0.008,-0.002,0.002,0.008,0.018\}.
\end{align*}
The older draft used $\mathcal V_a$ as held-out cases; they therefore cannot
serve as a newly untouched final test. Symmetry-generated orbit copies stay
in the same split. Training uses no interior reference-solution labels;
validation and test references still count toward setup cost.

Use float64 and smooth $\tanh$ networks so the third-order original residual
and optimization derivatives exist. A pilot architecture with two hidden
layers of width 24 has 697 parameters for input $(z,a)$ and 721 for
$(t,s,a)$; report these counts rather than claiming identical budgets.
Use fixed affine input scaling from the declared domains, including
$\widetilde z=4z-5$ and $\widetilde a=a/0.02$; compute physical residual
derivatives with all inverse scaling factors. Pair B/H initializations and
collocation seeds so their primary difference is the analytical prior.
Start with B, H and D, seeds 0 and 1, 128 residual points per training
parameter and 1000 Adam updates at learning rate $10^{-3}$. Sample the
direct arm in $(t,z)$ and transform to $(t,s)$ with the correct quadrature
weight. Draw independent collocation coordinates using the same declared
sampler. If the pilot passes the mathematical/reference and numerical
finiteness gates, confirm with seeds 0--4 at 64, 128 and 256 points per
parameter. Test H5 only as the separately declared ablation. Run sequentially
within the available CPU or 6 GB GPU budget and record the actual hardware.

This pilot budget is a starting protocol, not an assertion of adequate
capacity or convergence. Use $\mathcal V_a$ for any documented budget or
conditioning adjustment, freeze the final configurations before opening
$\mathcal S_a$, and retain all seeds and failures. Report both matched-step
comparisons and error versus measured time. The original-volume objective
for B/H can be evaluated using the factor $7/24$ above; normalized objectives
use the corresponding volume. Use the same fixed nondimensional residual
scales across arms, not method-dependent denominators chosen after training.
Distinct quotient states, original-domain samples, derivative evaluations
and optimizer updates are separate cost counts.

\subsection{Metrics, artifacts and conditions for a conclusion}
On independent grids and all held-out parameters report errors in
$F,F',F'',F'''$ and reconstructed $\phi$ derivatives; raw and fixed-scale
normalized reduced/original residuals; data errors; $Q_5$ or orbit dependence
for D; guard margins; seed medians/ranges; and failures. Include the reference
solver's error budget. A third derivative supplied directly by the ODE's
own normal-form right-hand side is not an independent residual certificate;
validate with refinement, solver agreement and an independent field/residual
evaluation route. Parameter-derivative accuracy must be measured separately
if claimed.

The preregistered primary target is at most $10^{-5}$ held-out absolute
profile error, with the reference discrepancy threshold already specified.
Failure to reach it remains a reportable outcome. To claim an improvement,
show the prior contrast H versus B and coordinate contrast B versus D
across seeds and budgets, with residual/derivative errors and costs, not
only the best profile-error seed. A formal error certificate additionally
requires the stability, reconstruction and continuous-residual bounds in
Proposition~\ref{prop:error-transfer-core}.

Store an equation/certificate hash, code version, configuration, seeds,
training histories, reference profiles, held-out predictions, errors,
guard checks and timing records. Separate analytical setup, reference
generation, training and full query cost, including reconstruction,
derivatives, verification and fallback. An amortized speedup requires
measuring all of these on the actual query task; historical break-even
arithmetic is not new evidence. Only after a continued subgroup or an
unreduced correction problem is verified should $\delta=m-3$ become an ML
input. A separately certified fixed-$m=3$ experiment may supply downstream
numerical evidence, but it cannot bypass or resolve the active near-cubic gate.


\clearpage
\appendix
\section{Historical numerical reports awaiting matching evidence}
\label{app:historical}
\textbf{Evidence status.} The numerical values below are preserved from the
earlier manuscript; they are not new experiments performed for this paper
update. The recovered archive does not contain the matching raw
Duffing/logistic/fusion and strict v0.2 family artifacts or the newer generic
backend needed to reproduce all of these claims. They must not support a
publication performance claim until the matching sources, configurations
and results are recovered and verified. In particular, the older comparison
between a hard-IVP PINN and a $P_2$ correction needs its implementation
explained: both can share the same hard-data polynomial.

Historical notation in this appendix maps $\sigma$ to the new quotient
coordinate $z$, $\eta_{\rm ic}$ to the new data parameter $a$, and $\rho=t$
to an alternative orbit coordinate; the canonical normalized orbit
coordinate in the new derivation is $y=\log(t/t_*)$.

\subsection{Reported ML baselines kept outside the Lie pipeline}
The ML-only executable imports no Lie-reduction code.  Its first experiment is
a bounded fixed-IVP reproduction of the Duffing comparison in the Taylor-model
PINN paper \cite{tmpinn},
\[
 \dot x=y,\qquad \dot y=x-x^3-\delta_Dy,
 \quad (x(0),y(0),\delta_D)=(-0.224,-0.269,0.327).
\]
Here $\delta_D$ is the Duffing damping coefficient and is unrelated to
$\delta=m-3$. All three methods use one hidden layer of width 64 and three
seeds. Table~\ref{tab:duffing} gives relative $L^2$ medians after 3000 Adam
steps; the bracket is the full seed range for TM-PINN.
\begin{table}[htbp]
\centering\small
\begin{tabular}{@{}ccccc@{}}
\toprule
Horizon & vanilla PINN & higher-order PINN & TM-PINN & TM range \\
\midrule
$T=1$ & $2.050\times10^{-4}$ & $1.914\times10^{-4}$ & $7.585\times10^{-5}$
 & $[4.491\times10^{-5},1.001\times10^{-4}]$ \\
$T=3$ & $6.555\times10^{-2}$ & $9.841\times10^{-2}$ & $1.756$
 & $[8.678\times10^{-2},8.894]$ \\
\bottomrule
\end{tabular}
\caption{Historical, not re-executed: bounded Duffing fixed-IVP reproduction. It tests the reported
short-time advantage and long-time degradation, not the paper's full
parameter-family solution map.}
\label{tab:duffing}
\end{table}

The second experiment follows the 2026 PINN--PIKAN logistic benchmark
\cite{pikan} with 100 uniform points, Adam step size $0.01$, 10000 iterations
and ten seeds. The MLP has 121 parameters; the PIKAN has 112 trainable
parameters after unused symbolic affine parameters are frozen. PIKAN's
relative-$L^2$ median/mean are $8.580\times10^{-5}/1.148\times10^{-3}$, while the
MLP values are $9.740\times10^{-1}/6.895\times10^{-1}$. Thus the reported direction
is reproduced, but the numerical values are not: several MLP seeds converge
to low-residual wrong branches. Median CPU training times are 63.71 s and
7.24 s, respectively. PirateNet \cite{piratenet} remains a literature
comparator because its JAX/Fourier/causal-training protocol was not reproduced.

\subsection{Pure Lie baseline}
The Lie-only executable imports no neural library. For the printed master
equation it verifies the supplied $m=3$ multiplier identity for
$X_5=t\partial_t-2s\partial_s$, integrates its flow, checks
$\sigma=st^2$ and $Q=-t\phi_t+2s\phi_s=0$, derives the third-order reduced
equation, and verifies
\[
 \E_3[F(st^2)]=t^{-1}\mathcal R_0[F],\qquad
 F_0=C+K\sigma^{1/3},\quad
 K^3=\frac{27g}{1-4g\beta_{\rm tr}}.
\]
The same supplied identity at $m=4$ returns \texttt{unsupported}. This is a
negative control for $X_5$, not a proof that no other $m=4$ symmetry exists.

\subsection{Certificate-gated Lie--ML prototype}
The fusion executable refuses an unverified Lie-only certificate. It fixes
$g=1$, $\beta_{\rm tr}=0.1$, $m=3$, $\sigma\in[1,1.5]$, and changes only
the initial curvature by
$F''(1)\mapsto(1+\eta_{\rm ic})F_0''(1)$ with $\eta_{\rm ic}=0.01$;
$\delta=m-3$ remains zero. A coarse/medium/fine DOP853 tolerance ladder uses
92, 197 and 377 function evaluations; its first two differences from fine are
$2.025\times10^{-11}$ and $2.065\times10^{-14}$. Fine DOP853 and Radau agree
to $3.281\times10^{-13}$.
The three neural methods have 673 parameters, use 1000 Adam steps and retain
all three seeds. Table~\ref{tab:fusion} reports seed medians.
\begin{table}[htbp]
\centering\small
\begin{tabular}{@{}lrrrr@{}}
\toprule
Method & $\max|F-F_{\rm ref}|$ & raw $\mathcal R_0$ & normalized
 $\mathcal R_0$ & original $\E_3$ \\
\midrule
hard-IVP PINN & $3.425\times10^{-5}$ & $4.701\times10^{-1}$ & $4.496\times10^{-2}$ & $4.711\times10^{-1}$ \\
$P_2+h^3N_\theta$ & $2.670\times10^{-6}$ & $1.804\times10^{-1}$ & $2.535\times10^{-2}$ & $1.831\times10^{-1}$ \\
$P_4+h^5N_\theta$ & $3.467\times10^{-5}$ & $5.354\times10^{-2}$ & $5.470\times10^{-3}$ & $5.458\times10^{-2}$ \\
\bottomrule
\end{tabular}
\caption{Historical, not re-executed: scalar-benchmark prototype; $h=\sigma-1$. Profile and residual rankings disagree; both
are therefore retained. Direct and lifted residuals agree to about $10^{-14}$.}
\label{tab:fusion}
\end{table}

The $P_2$ correction has the smallest profile error, whereas the $P_4$
correction has the smallest residual. Fine DOP853 takes about 0.0034 s for this
single IVP, versus 3.6--4.2 s training per neural seed. The prototype therefore
does not establish an ML speedup. Cubic PIKAN is not transferred to this
third-order residual because its third-derivative smoothness has not passed a
separate gate.

\subsection{Stage goal: a data-compatible parameter-family certificate}
The bounded v0.2 goal is not the unresolved generic-$m$ classification. It is
to close one local, falsifiable chain on the fixed cubic branch:
\[
 (\E_3,X_5,\mathcal D_\eta)
 \longrightarrow (\mathcal R_0,\mathcal C_{0.2})
 \longrightarrow \widehat F_\theta(\sigma,\eta)
 \longrightarrow \E_3[\widehat F_\theta(st^2)].
\]
Here $\rho=t$, $\sigma=st^2$, and $\mathcal C_{0.2}$ is accepted only if the
chart rank, orbit speed, fibre rank, exact pullback, and data checks all pass.
On the invariant hypersurface $\Sigma_0=\{st^2=1\}$ we prescribe
\[
 \phi-K=0,\qquad t^{-2}\phi_s-\frac K3=0,\qquad
 t^{-4}\phi_{ss}+\frac{2K}{9}(1+\eta_{\rm ic})=0,
 \quad K=\sqrt[3]{45}.
\]
The prolonged action annihilates the surface residual and all three data
residuals off shell, and their pullbacks contain no $\rho$. Thus they become
\[
 F(1)=K,\qquad F'(1)=K/3,\qquad
 F''(1)=-\frac{2K}{9}(1+\eta_{\rm ic}).
\]
This is a strict sufficient compatibility test for this constructed data
family, not a statement that arbitrary canal boundaries are symmetry
compatible. The certificate excludes $\rho=0$, $F'=0$, and
\[
 a(\sigma,F')=-\frac{2g\sigma}{(F')^2}
 +8(1-\beta_{\rm tr}g)\sigma^3F'=0;
\]
it records these as separate singular exits rather than dividing them away.

For $\eta_{\rm ic}\in[-0.02,0.02]$ the conditional profile is
\[
 \widehat F_\theta(\sigma,\eta_{\rm ic})
 =P_2(\sigma,\eta_{\rm ic})
  +(\sigma-1)^3N_\theta(\sigma,\eta_{\rm ic}),
\]
so the three data jets are exact for training and unseen parameters. Training
and held-out parameter sets are respectively
\[
\begin{aligned}
 \mathcal T_\eta&=\{-0.02,-0.015,-0.01,0,0.01,0.015,0.02\},\\
 \mathcal H_\eta&=\{-0.0125,-0.005,0.005,0.0125\}.
\end{aligned}
\]
They are disjoint, and every run retains seeds $0,1,2$. All eleven numerical
references are sampled regular and DOP853/Radau differ by at most
$3.312\times10^{-13}$. All twelve held-out evaluations complete; their
maximum profile errors range from $2.636\times10^{-5}$ to
$1.670\times10^{-4}$, compared with about $2.05\times10^{-2}$ for $P_2$ and
$2.4\times10^{-3}$ for $P_4$. The initial-jet error is at most
$1.111\times10^{-16}$ and direct-versus-lifted PDE residuals differ by at most
$1.616\times10^{-14}$.

Median training, DOP853-query, and neural-query times are respectively
$4.625$ s, $0.00326$ s, and $0.00262$ s. Their measured arithmetic break-even
is about 7170 repeated queries, before deployment and uncertainty costs.
Consequently this experiment supports held-out interpolation beyond finite
Taylor truncations on the declared regular family, but still does not prove an
ML speedup, an interval error bound, or a solution on a singular branch.

\section{Historical implementation record}
The earlier manuscript recorded the following four-day sprint. This is a
historical work record subject to the evidence status in Appendix~\ref{app:historical}.
\begin{center}
\small
\begin{tabularx}{\linewidth}{@{}lXX@{}}
\toprule
Day & Work & Result \\
\midrule
1 & Lock notation and finish the Lie-only reduction certificate. &
$m=3$ verified; the supplied $X_5$ identity at $m=4$ rejected. \\
2 & Run the TM-PINN and PINN--PIKAN literature baselines. &
All declared seeds retained; directional and negative results separated. \\
3 & Run the certificate-gated SVE fusion against DOP853 and Radau. &
Three methods and three seeds completed; lifting checked in the original PDE. \\
4 & Consolidate JSON, CSV, figures, manuscript and interim report. &
Reproducible commands, evidence limits and next gates recorded. \\
\bottomrule
\end{tabularx}
\end{center}
Having tested one constructed initial-curvature family, the next gates are to
select a subgroup compatible with physically sourced initial/boundary data,
analyze the recorded rank and leading-coefficient singular exits, and establish
a smoothness gate before any KAN transfer to a third-order residual. The case
$m=3+\delta$ remains separate until its symmetry-persistence calculation. A
local series, rational continuation, or matched expansion should be selected
only after analyzing the reduced equation's singularities and admissible data.


\section{Example implementation: a declared SVE candidate ansatz}
\label{app:sve-listing}
This worked listing belongs to the scalar example, not the general theorem.
Its finite trial space is explicitly bounded; the shorter independent
checks used for this paper update verify supplied identities, not maximality.

The following compact jet-space implementation uses SymPy's exact
differentiation, polynomial collection, and matrix nullspace operations. It
spells out the prolongation instead of relying on a black-box Lie solver;
the documented SymPy PDE interface includes utilities such as separation,
which are different from a general PDE symmetry classifier \cite{sympydocs}.
The code below computes candidates in a declared polynomial ansatz for the
regular $m=3$ branch. It does not compute the generic-$m$ algebra: that branch
needs a separate determining-system run and a branch-aware split of the
noninteger powers on $phi_t>0$. Increase the ansatz only after examining the
determining system; completeness requires solving that full linear
differential system, including its arbitrary-function branches, not merely
enlarging a finite polynomial trial space.

\begin{lstlisting}[language=Python]
import sympy as sp

t, s, m = sp.symbols("t s m", real=True)
g, Ag, xi_sed = sp.symbols("g Ag xi_sed", positive=True)
R = 3                         # order of the SVE master PDE
jets = {(i, j): sp.Symbol(f"u{i}{j}", real=True)
        for i in range(R + 2) for j in range(R + 2 - i)}
u = jets[(0, 0)]

def D(expr, axis):            # axis 0=t, 1=s; total derivative
    out = sp.diff(expr, (t, s)[axis])
    for (i, j), uj in jets.items():
        if i + j <= R:
            nxt = (i + (axis == 0), j + (axis == 1))
            out += sp.diff(expr, uj) * jets[nxt]
    return sp.expand(out)

Dt, Ds = lambda f: D(f, 0), lambda f: D(f, 1)
us, ut = jets[(0, 1)], jets[(1, 0)]
utt, uss, uts = jets[(2, 0)], jets[(0, 2)], jets[(1, 1)]
beta = m * Ag * xi_sed
E = Dt(us*utt - g*uss/us**2) - g*beta*Ds(ut**(m-1)*uts/us)
E3 = sp.cancel(E.subs(m, 3))

def prolong(tau, xi, eta, r=R):
    e = {(0, 0): eta}
    for n in range(1, r + 1):
        for i in range(n + 1):
            j = n - i
            if i: parent, axis = (i - 1, j), 0
            else: parent, axis = (i, j - 1), 1
            a, b = parent
            e[(i, j)] = sp.expand(
                D(e[parent], axis)
                - jets[(a + 1, b)] * D(tau, axis)
                - jets[(a, b + 1)] * D(xi, axis))
    return e

def prX(F, tau, xi, eta, r=R):
    e = prolong(tau, xi, eta, r)
    ans = tau*sp.diff(F, t) + xi*sp.diff(F, s) + eta*sp.diff(F, u)
    for (i, j), ej in e.items():
        if i + j and (i + j) <= r:
            ans += ej * sp.diff(F, jets[(i, j)])
    return sp.expand(ans)

# Candidate search only: total-degree <= 2 in (t,s,u).
mon = [t**a*s**b*u**c for a in range(3) for b in range(3-a)
       for c in range(3-a-b)]
def ansatz(prefix):
    cs = sp.symbols(f"{prefix}0:{len(mon)}")
    return sum(c*p for c, p in zip(cs, mon)), cs

tau, c_tau = ansatz("ct")
xi, c_xi = ansatz("cs")
eta, c_eta = ansatz("cu")
cvars = c_tau + c_xi + c_eta
lead = jets[(3, 0)]       # regular branch: solve E3 for phi_ttt
lead_rule = sp.solve(E3, lead, dict=True)[0]
determining = sp.cancel(prX(E3, tau, xi, eta).subs(lead_rule))
num = sp.together(determining).as_numer_denom()[0]
free_jets = [jets[J] for J in jets
             if 0 < sum(J) <= R and J != (3, 0)]
P = sp.Poly(num, *(free_jets + [t, s, u]), domain="EX")
eqs = P.coeffs()                       # split all free-jet monomials
M, rhs = sp.linear_eq_to_matrix(eqs, cvars)
basis = M.nullspace()                  # exact candidate generators
generators = []
for v in basis:
    gen = tuple(sp.expand(f.subs(dict(zip(cvars, v))))
                for f in (tau, xi, eta))
    test = sp.cancel(prX(E3, *gen).subs(lead_rule))
    assert sp.simplify(test) == 0       # verify each candidate on E3=0
    generators.append(gen)
check_x5 = sp.cancel(prX(E3, t, -2*s, 0).subs(lead_rule))
assert sp.simplify(check_x5) == 0       # independent m=3 anchor check
\end{lstlisting}

Read the steps as an executable protocol. First form the PDE using independent
jet symbols, not unevaluated derivatives of an unknown function. The total
derivative operator raises each jet index and the recursion constructs
$\eta_J=D_i\eta_{J-e_i}-u_{J-e_i,t}D_i\tau-u_{J-e_i,s}D_i\xi$.
Second, solve the regular PDE branch for a chosen leading jet and substitute
it into $\operatorname{pr}^{(3)}X(E)$. For the SVE equation the leading
coefficient is nonzero on $\phi_s>0$; record that condition and analyze
singular branches separately. Third, clear only denominators known nonzero on
the branch, split with respect to all remaining jets (and then the base
coordinates), and solve the resulting homogeneous linear system. The matrix
nullspace returns a basis within the stated ansatz. For a system, repeat the
criterion for every $E_A$ and eliminate a compatible set of leading jets;
check differential consequences whenever the system is not in a regular
solved form.

For generic symbolic $m$, the factors $\phi_t^{m+k}$ are not ordinary
polynomial monomials. Split them as distinct smooth functions on the chosen
$\phi_t>0$ branch, or analyze rational/integer $m$ strata separately; then
solve $m=3$ again without first dividing by $m-3$. For
$m=3+\delta$, expand powers as
$(\phi_t/U_*)^{2+\delta}=(\phi_t/U_*)^2(1+\delta\log(\phi_t/U_*))+O(\delta^2)$ and retain the
logarithmic basis element in the split. To certify maximality, replace the
finite polynomial ansatz by unknown functions $\tau(t,s,\phi)$,
$\xi(t,s,\phi)$, and $\eta(t,s,\phi)$, integrate the determining PDEs (or
complete them to an involutive differential system), branch on every
parameter/rank condition, and verify the final basis by substitution. Exact
differential completion is the completeness gate; a numerical SVD or finite
degree ansatz is a discovery aid \cite{approxalg}.

For every basis element, run an independent invariance check by substituting
the solved-jet rules into $\operatorname{pr}^{(3)}X(E)$ and simplifying to
zero. For coefficient triples $(\tau_i,\xi_i,\eta_i)$, compute
\[
 [X_i,X_j]\ \longleftrightarrow\
 X_i(\tau_j,\xi_j,\eta_j)-X_j(\tau_i,\xi_i,\eta_i).
\]
Express these triples in the basis and check the structure constants and Jacobi.
The adjoint
representation then supports the optimal-system calculation. For $X_5$,
verify $X_5(st^2)=0$ and $X_5(\phi)=0$, substitute $\phi=F(st^2)$ directly,
and simplify $t\E_3-\mathcal R_0[F]$ to zero. Finally exponentiate candidate
generators and check the finite transformation on the physical domain; local
infinitesimal validity alone does not establish a globally admissible action.
